\chapter{Discussion}
\label{chap:Discussion}

The results obtained in this project show that deep learning models can successfully learn spatial patterns from meteorological data and generate short-term weather forecasts. The U-Net architecture used in this study demonstrated the ability to capture large-scale temperature structures and produce reasonable predictions for short forecast horizons. However, several aspects of the approach could be further explored to improve performance and better evaluate the effectiveness of the model.
First, although the experiments mainly focused on a U-Net encoder–decoder architecture, other neural network architectures could also be investigated. For example, recurrent models such as LSTM or GRU networks could be used to better capture temporal dependencies in atmospheric data. Hybrid architectures such as ConvLSTM or transformer-based models could also be explored, as they combine spatial and temporal modeling capabilities. Testing several model families would provide a clearer understanding of which architectures are most suitable for short-term weather prediction.
Another important aspect concerns the computational efficiency of the models. Future experiments could compare the training and inference times of different architectures. It is possible that lighter models with fewer parameters could achieve comparable prediction accuracy while requiring significantly less computational resources. Such comparisons would be particularly important for operational forecasting systems, where both accuracy and computational cost are critical factors.
The computational experiments could also be extended by running the models on more powerful hardware. In this project, training was performed using GPUs available on Google Colab and a local NVIDIA GTX 1660 Ti. Running the training pipeline on more powerful accelerators, such as an NVIDIA A100 GPU, could allow the use of larger batch sizes, longer training runs, and larger datasets. This would make it possible to explore the effect of increasing the number of training epochs and determine whether additional training leads to improved accuracy. However, this would require careful monitoring in order to avoid overfitting, as excessive training can lead the model to memorize the training data rather than generalize to unseen observations.
Another possible improvement concerns the design of the model inputs. Instead of relying only on the temperature values at a single location or grid cell, the model could incorporate temperature information from geographically neighboring regions. Since atmospheric processes are spatially correlated, including surrounding temperature values could help the model better capture the evolution of local weather conditions and improve prediction accuracy.
Further improvements could also involve modifying the hyperparameters of the U-Net architecture. For example, the number of encoder and decoder layers, the number of convolutional filters, or the types of layers used in the network could be adjusted to study their impact on model performance. Systematic hyperparameter tuning would allow a more comprehensive exploration of the architecture and potentially lead to better forecasting results.
Another limitation of the current study is that only one meteorological variable, surface temperature, was used as input. In reality, atmospheric dynamics depend on multiple interacting variables such as wind speed, humidity, precipitation, and pressure. Future experiments could therefore incorporate additional meteorological variables into the model. For instance, variables such as rainfall or wind fields could provide additional information about atmospheric dynamics. It may also be possible to design a multi-variable model capable of predicting several atmospheric variables simultaneously, which could better reflect the complexity of real weather systems.
In addition, the geographical scope of the experiments could be expanded. The current study focuses on a specific region over Europe, but it would be interesting to extend the experiments to other parts of the world. Training models on larger geographical regions or even global datasets could provide a better understanding of how well deep learning models generalize across different climates and atmospheric regimes. However, expanding the spatial domain would significantly increase the size of the datasets and the computational requirements for model training. As a result, careful consideration would need to be given to the available computational resources.
Finally, it would be useful to compare the performance of the implemented model with more advanced AI-based weather forecasting systems. Recent models such as GraphCast, FourCastNet, or Pangu-Weather have demonstrated strong forecasting capabilities using large-scale datasets and advanced architectures. Although reproducing such models is beyond the scope of this project, comparing their reported performance with the results obtained here could help evaluate the relative efficiency of the proposed approach and identify potential directions for improvement.
Overall, while the results obtained in this project are encouraging, further experimentation with different architectures, input configurations, meteorological variables, and computational setups would be necessary to fully assess the potential of deep learning methods for weather forecasting.
